\section{Parabolics with a nontrivial chamber stabilizer}
We now allow flag-transitive coset geometries whose acting groups are not
faithful. The type sets remain finite and nonempty. In a geometry
$\alpha=(A,(A_x)_{x\in S})$ put $A_E=\bigcap_{x\in E}A_x$, with
$A_\varnothing=A$, and call $A_S$ the Borel subgroup. It need not be
trivial. The residue of the standard flag of types $E$ is the coset
geometry of $A_E$ with parabolics $A_{E\cup\{x\}}$, $x\notin E$.
To see this identification, take an element incident with that flag;
flag transitivity supplies a common representative in $A_E$. Two such
representatives define the same element exactly when their quotient
lies in $A_E\cap A_x$. This also identifies incidence by intersecting
the corresponding cosets.

For a coset geometry of rank at least two the incidence graph is
connected if and only if its maximal parabolics generate the acting
group. Indeed, if $H$ is that generated subgroup, paths from the
standard flag remain among cosets with representatives in $H$.
Conversely multiplication successively by elements of the parabolics
provides paths through incident cosets; when two successive steps have
the same type, an intervening different type gives the same conclusion.
Thus the components are the translates of this component. Applied to
every residue, this proves that residual connectedness is equivalent to
\begin{equation}\label{eq:rc-borel}
 A_E=\langle A_{E\cup\{x\}}:x\notin E\rangle
 \quad\text{if }|S\setminus E|\geq2.
\end{equation}
For rank one residual connectedness imposes no condition.

\begin{lemma}[Generation of parabolics]\label{lem:parabolic-join}
If $\alpha$ is residually connected, then
\[
 A_E=\langle A_S,\ A_{S\setminus\{x\}}:x\notin E\rangle.
\]
Consequently, for any nonempty family $(E_\nu)$ of subsets of $S$,
\begin{equation}\label{eq:parabolic-join}
 \langle A_{E_\nu}:\nu\rangle=A_{\bigcap_\nu E_\nu}.
\end{equation}
\end{lemma}
\begin{proof}
Induct on $|S\setminus E|$. For zero this is $A_S=A_S$; for one
the displayed generator is $A_E$ itself. For larger complements,
use (\ref{eq:rc-borel}) and the induction hypothesis on each child.
Every missing type remains missing in at least one child, so all and
only the indicated minimal parabolics occur. For the second assertion,
each group on the left lies in the right-hand group. Conversely, every
minimal parabolic required to generate the right-hand group occurs
in some $A_{E_\nu}$, and $A_S$ lies in every one. The nonempty-family
hypothesis is necessary when $A_S\ne1$.
\end{proof}
The explicit $A_S$ in the first formula must not be dropped at
$E=S$. This avoids the empty-generated-subgroup problem in the
source's background equivalences.

\section{A general orbit construction}
Let $\alpha=(A,(A_x)_{x\in S})$ and
$\beta=(B,(B_i)_{i\in I})$ be flag-transitive coset geometries.
Suppose a specified homomorphism $B\to\Aut(A)$ and an action of $B$
on the indexed set $S$ satisfy $b(A_x)=A_{bx}$.
The indexed action is part of the data, even if different maximal
parabolics happen to coincide. Let $\mathcal K$ be the set of orbits
on $S$, choose $x_L\in L$ for each $L\in\mathcal K$, and assume
\begin{equation}\label{eq:general-orbits}
 B_Hx_L\cap B_Kx_L=B_{H\cup K}x_L
 \quad(H,K\subseteq I,\ L\in\mathcal K).
\end{equation}
For $T=I\sqcup\mathcal K$, define subgroups of $G=A\rtimes B$ by
\begin{equation}\label{eq:general-maximal}
 G_i=A_{S\setminus\bigcup_L B_ix_L}\rtimes B_i\quad(i\in I),
 \qquad G_L=A_L\rtimes B\quad(L\in\mathcal K).
\end{equation}
Every indicated subset of $S$ is invariant under its second factor,
so these are indeed subgroups. For $J\subseteq T$ write
$J_B=J\cap I$, $J_A=J\cap\mathcal K$, and set
\[
 E(J)=S\setminus\bigcup_{L\in\mathcal K\setminus J_A}B_{J_B}x_L.
\]

\begin{theorem}[Orbit construction with a Borel subgroup]\label{thm:general-twist}
The coset system in (\ref{eq:general-maximal}) is a flag-transitive
geometry. All its standard parabolics, including its Borel subgroup, are
\begin{equation}\label{eq:general-parabolics}
 G_J=A_{E(J)}\rtimes B_{J_B},\qquad G_T=A_S\rtimes B_I.
\end{equation}
It is finite if and only if both input geometries are finite. It is
residually connected if both inputs are residually connected, and firm
if both inputs are firm. Its numbers of chambers through panels are
\begin{align}
 [G_{T\setminus\{i\}}:G_T]&=[B_{I\setminus\{i\}}:B_I]
       &&(i\in I),\label{eq:old-panel}\\
 [G_{T\setminus\{L\}}:G_T]&=[A_{S\setminus B_Ix_L}:A_S]
       &&(L\in\mathcal K).\label{eq:new-panel}
\end{align}
In particular, if both inputs are thin and $B_I$ fixes each chosen
$x_L$, the new geometry is thin. Faithfulness of either input is not
needed for this sufficient condition.
\end{theorem}
\begin{proof}
Intersections of subgroups $A_E\rtimes B_H$ are computed coordinate by
coordinate. The intersection of their $A$-parts is obtained by taking
the union of their subscripts. Different $B$-orbits are disjoint, and
iteration of (\ref{eq:general-orbits}) says that the intersection of
the old supports is the support for $B_{J_B}$. Intersecting with the
orbit parabolics deletes precisely the orbits in $J_A$. Taking
complements proves (\ref{eq:general-parabolics}). It uses intersections
of maximal parabolics, not an unjustified formula for generated supports.

To prove flag transitivity, project a flag of $G$-cosets to $B$.
Old types give pairwise-intersecting maximal-parabolic cosets; orbit
types project to all of $B$. Coset intersection from the first section
gives a common $b\in B$. After left translation by $(1,b^{-1})$,
each coset has a representative $(a,1)$. Such a coset is the Cartesian
product of a coset of a standard parabolic of $A$ and the corresponding
subgroup of $B$. Pairwise incidence therefore gives a pairwise-intersecting
family of standard-parabolic $A$-cosets. Their common representative,
again by the coset intersection lemma, realizes the entire flag.
Thus every flag is an image of a standard flag and extends to a chamber.
This step by itself does not require the orbit identity; that identity
is needed for the exact parabolic description and the later conclusions.

Suppose both inputs are residually connected, and fix $J$ with at
least two missing types. We prove generation of $G_J$ by the next
parabolics. If no orbit type is missing, all their $A$-parts are $A_S$;
(\ref{eq:rc-borel}) for $B$ gives the required generation. If no old
type is missing, choose two distinct missing orbits $L,M$. The two
corresponding next parabolics both contain $B_I$. Their $A$-subscripts
are $E(J)\cup L$ and $E(J)\cup M$, whose intersection is $E(J)$.
Lemma~\ref{lem:parabolic-join} gives the entire $A$-part.

In the remaining case choose a missing old type $i$ and a missing
orbit type $L$. The parabolic $G_{J\cup\{L\}}$ contains the full
$B_{J_B}$. Write $E'=E(J\cup\{i\})$. The group generated by these
two next parabolics therefore contains $A_{bE'}$ for every
$b\in B_{J_B}$. In each remaining orbit, $B_{J_B\cup\{i\}}x_M$
is nonempty and contains $x_M$. Its $B_{J_B}$-saturation is exactly
$B_{J_B}x_M$. Consequently
\[
 \bigcap_{b\in B_{J_B}}bE'=E(J).
\]
There are only finitely many distinct subsets of the finite set $S$,
so Lemma~\ref{lem:parabolic-join} applies directly to these conjugate
parabolics. They generate $A_{E(J)}$, proving the remaining case.
The residue identification and connectivity criterion already proved
now establish residual connectedness in every residue.

Putting $J=T\setminus\{t\}$ in
(\ref{eq:general-parabolics}) gives the two panel formulas.
For an old type the $A$-parts agree; for an orbit type the $B$-parts
agree. These observations prove the stated indices even when the
Borel subgroups are nontrivial. If the inputs are firm, the old indices
are at least two. The orbit $B_Ix_L$ contains $x_L$, so
$A_{S\setminus\{x_L\}}\leq A_{S\setminus B_Ix_L}$; its index over
$A_S$ is at least two by firmness of $\alpha$. This proves firmness.
If the inputs are thin and $B_Ix_L=\{x_L\}$, all the displayed
indices are two, proving the sufficient thinness assertion.

Finally the number of chambers is the index of $G_T$. The elementary
coset index formula gives
\[
 [G:G_T]=[A:A_S][B:B_I],
\]
with the equality interpreted as a product of cardinalities.
For a finite number of types, a geometry has finitely many elements
if and only if it has finitely many chambers: every element extends
to a chamber and a chamber is a tuple of one element of each type.
The same argument applies to each input. Thus the product is finite
exactly when both input geometries are finite. No finiteness of the
possibly unfaithful groups $A$ and $B$ was assumed.
\end{proof}

\begin{corollary}[Faithful regular inputs]\label{thm:twist}
For faithful regular hypertopes $\alpha,\beta$, the construction is a
regular hypertope with full type-preserving automorphism group $A\rtimes B$.
Writing $A(R)=\langle a_x:x\in R\rangle$, its parabolics are
\begin{equation}\label{eq:regular-support}
 G_J=A\!\left(\bigcup_{L\notin J_A}B_{J_B}x_L\right)\rtimes B_{J_B}.
\end{equation}
Its distinguished involutions are the $b_i$ and the $a_{x_L}$.
\end{corollary}
\begin{proof}
The faithful regular model has $A_S=B_I=1$ and
$A_E=A(S\setminus E)$, as proved in the first section. The preceding
theorem supplies flag transitivity, residual connectedness and thinness.
Its Borel subgroup is trivial, and the panel groups are precisely the
cyclic groups on the indicated involutions. Chamber propagation shows
that the displayed regular action is the full type-preserving group.
\end{proof}
This recovers the general finite-type twisting theorem and its regular
specialization in Piedade--Tranchida [2, Theorem 3.16 and Corollary 3.17].
The exact panel formula also explains why passing to a faithful quotient
of $B$ cannot silently preserve an arbitrary action on the types of
$\alpha$: its Borel subgroup must act trivially for that action to descend.
